\documentclass[aspectratio=169]{beamer}

% ===================== TEMA =====================
\usetheme{metropolis}
\usepackage[utf8]{inputenc}
\usepackage[T1]{fontenc}
% es-noshorthands: sin esto babel convierte "o en "º" dentro del código
\usepackage[spanish,es-tabla,es-noshorthands]{babel}
\usepackage{graphicx}
\usepackage{booktabs}
\usepackage{tikz}
\usepackage{listings}
\usepackage{qrcode}
\usepackage{pifont}
\usetikzlibrary{positioning,arrows.meta,calc,fit,shapes.misc,decorations.pathreplacing}

% Color de acento propio, en línea con el apunte
\definecolor{accentblue}{RGB}{30,70,140}
\setbeamercolor{progress bar}{fg=accentblue, bg=accentblue!20}
\setbeamercolor{frametitle}{bg=accentblue}
\setbeamercolor{alerted text}{fg=accentblue}
\setbeamertemplate{frame numbering}[fraction]

% ===================== CÓDIGO =====================
\definecolor{codebg}{RGB}{245,246,248}
\definecolor{codecomment}{RGB}{110,120,130}
\definecolor{codestring}{RGB}{160,60,20}
\lstset{
    basicstyle=\ttfamily\scriptsize,
    backgroundcolor=\color{codebg},
    keywordstyle=\color{accentblue}\bfseries,
    commentstyle=\color{codecomment}\itshape,
    stringstyle=\color{codestring},
    showstringspaces=false,
    upquote=true,
    frame=single, rulecolor=\color{accentblue!30},
    framesep=3pt, xleftmargin=4pt, xrightmargin=4pt,
    columns=fullflexible, keepspaces=true,
    literate={á}{{\'a}}1 {é}{{\'e}}1 {í}{{\'i}}1 {ó}{{\'o}}1 {ú}{{\'u}}1 {ñ}{{\~n}}1,
}
\lstdefinelanguage{dts}{
    morekeywords={compatible,reg,status,gpios,aliases,chosen,label,pinctrl-0,pinctrl-names,clock-frequency},
    sensitive=true,
    morecomment=[l]{//},
    morecomment=[s]{/*}{*/},
    morestring=[b]",
}
\lstdefinelanguage{kconf}{morecomment=[l]{\#}, morekeywords={y,n}}

\newcommand{\si}{\textcolor{green!50!black}{\ding{51}}}
\newcommand{\no}{\textcolor{gray!60}{\ding{55}}}

% Repositorio del workspace
\newcommand{\repourl}{https://github.com/utn-fra-lse/cie2026-zephyr}

\title{Zephyr RTOS: una aplicación, muchas placas}
\subtitle{Devicetree, drivers estandarizados y portabilidad real}
\author{Fabrizio Carlassara\\
\footnotesize Laboratorio de Sistemas Embebidos}
\institute{UTN -- Facultad Regional Avellaneda}
\date{Congreso de Ingeniería Electrónica 2026}

\begin{document}

% =====================================================
\maketitle

% =====================================================
\begin{frame}{Hoja de ruta}
\begin{enumerate}
    \item El problema: código atado al fabricante
    \item Zephyr en cinco minutos
    \item \textbf{Devicetree}: describir el hardware en vez de programarlo
    \item \textbf{Modelo de drivers}: una API, muchas implementaciones
    \item \alert{Demo}: las mismas aplicaciones en placas de distintos fabricantes y arquitecturas
    \item Herramientas, límites y cómo empezar
\end{enumerate}
\end{frame}

% =====================================================
\section{El problema}

% =====================================================
\begin{frame}[fragile]{Prender un LED, cuatro fabricantes}
\begin{columns}[T]
\column{0.5\textwidth}
\textbf{\footnotesize STM32 (STM32Cube HAL)}
\begin{lstlisting}[language=C]
HAL_GPIO_WritePin(GPIOC, GPIO_PIN_13,
                  GPIO_PIN_SET);
\end{lstlisting}
\textbf{\footnotesize NXP (MCUXpresso SDK)}
\begin{lstlisting}[language=C]
GPIO_PinWrite(GPIO, 1, 4, 1U);
\end{lstlisting}
\column{0.5\textwidth}
\textbf{\footnotesize Espressif (ESP-IDF)}
\begin{lstlisting}[language=C]
gpio_set_level(GPIO_NUM_2, 1);
\end{lstlisting}
\textbf{\footnotesize Raspberry Pi (Pico SDK)}
\begin{lstlisting}[language=C]
gpio_put(25, true);
\end{lstlisting}
\end{columns}
\vspace{0.4em}
\centering
Cuatro APIs, cuatro sistemas de build, cuatro formas de configurar relojes y pines.\\
\alert{El pin y el periférico quedan \emph{escritos en el código}.}
\end{frame}

% =====================================================
\begin{frame}{El costo de casarse con un fabricante}
\begin{itemize}
    \item \textbf{Faltantes de componentes} (2021--2023): rediseñar la placa \emph{y} reescribir el firmware
    \item \textbf{Cambio de requerimientos}: más RAM, BLE, menor consumo\ldots{} $\Rightarrow$ otro micro
    \item \textbf{Equipos y formación}: cada familia es un mundo nuevo
    \item \textbf{Prototipo vs.\ producto}: se prueba en una placa de desarrollo y se fabrica con otra
\end{itemize}
\vspace{0.5em}
\begin{center}
\begin{tikzpicture}[box/.style={draw, rounded corners, minimum height=0.8cm, align=center, font=\small}]
\node[box, fill=green!15, minimum width=3.2cm] (app) {Lógica de la\\aplicación};
\node[box, fill=orange!20, right=0.2cm of app, minimum width=3.2cm] (hal) {HAL del\\fabricante};
\node[box, fill=gray!20, right=0.2cm of hal, minimum width=3.2cm] (hw) {Hardware};
\draw[decorate, decoration={brace, amplitude=6pt, mirror}, thick, red!70!black]
    ($(app.south west)+(0,-0.1)$) -- ($(hal.south east)+(0,-0.1)$)
    node[midway, below=8pt, font=\footnotesize, text=red!70!black] {mezclados en los mismos archivos};
\end{tikzpicture}
\end{center}
\end{frame}

% =====================================================
\begin{frame}[standout]
¿Y si la aplicación\\
\textbf{\emph{no supiera}} en qué placa corre?
\end{frame}

% =====================================================
\section{Zephyr en cinco minutos}

% =====================================================
\begin{frame}{¿Qué es Zephyr?}
\begin{columns}[T]
\column{0.58\textwidth}
\begin{itemize}
    \item RTOS de código abierto, proyecto de la \textbf{Linux Foundation} (2016)
    \item Licencia \textbf{Apache 2.0}: sin regalías, apto para productos comerciales
    \item Gobernanza abierta: Intel, Nordic, NXP, ST, Espressif, Google, Meta\ldots
    \item No es solo un kernel: \textbf{drivers, pilas de comunicación, sistema de build y herramientas}
    \item Versiones cada 6 meses, LTS con soporte extendido
\end{itemize}
\column{0.38\textwidth}
\centering
\renewcommand{\arraystretch}{1.3}
\begin{tabular}{rl}
\toprule
\multicolumn{2}{c}{\textbf{Zephyr v4.5.0}} \\
\midrule
\alert{\Large +1100} & placas \\
\alert{\Large 185} & fabricantes de placas \\
\alert{\Large 13} & arquitecturas \\
\alert{\Large +240} & drivers de sensores \\
\bottomrule
\end{tabular}

\vspace{0.3em}
{\tiny Contado sobre el árbol de v4.5.0}
\end{columns}
\end{frame}

% =====================================================
\begin{frame}{Arquitectura: dónde vive cada cosa}
\centering
\resizebox{0.92\textwidth}{!}{%
\begin{tikzpicture}[
    layer/.style={draw, text width=10cm, minimum height=0.9cm, align=center, font=\small},
    side/.style={draw, rounded corners, text width=3.2cm, align=center, font=\small, fill=yellow!20},
]
\node[layer, fill=gray!25] (hw) {Hardware\\\footnotesize RP2040 · STM32 · ESP32-C3 · nRF5340 · K64F · i.MX RT · \ldots};
\node[layer, fill=orange!20, above=0cm of hw] (hal) {HALs de los fabricantes \footnotesize(módulos, ocultos)};
\node[layer, fill=blue!20, above=0cm of hal] (drv) {Drivers \footnotesize(una implementación por periférico)};
\node[layer, fill=blue!10, above=0cm of drv] (api) {\textbf{APIs estándar}: GPIO, I\textsuperscript{2}C, SPI, UART, ADC, PWM, sensor\ldots};
\node[layer, fill=blue!15, above=0cm of api] (krn) {Kernel: threads, scheduler, semáforos, colas, timers};
\node[layer, fill=green!15, above=0cm of krn] (app) {\textbf{Aplicación} \footnotesize(portable)};

\node[side, anchor=west, minimum height=2.4cm] (dt) at ($(drv.east)!0.5!(api.east)+(1.2,0)$) {\textbf{Devicetree}\\\footnotesize ¿qué hardware hay\\\footnotesize y cómo está conectado?};
\node[side, anchor=west, minimum height=1.6cm] (kc) at ($(krn.east)!0.5!(app.east)+(1.2,0)$) {\textbf{Kconfig}\\\footnotesize ¿qué software\\\footnotesize se incluye?};
\draw[-{Latex}, thick, accentblue] (dt.west) -- (drv.east |- dt.west);
\draw[-{Latex}, thick, accentblue] (kc.west) -- (krn.east |- kc.west);
\node[font=\footnotesize\itshape, below=0.15cm of dt] {se resuelven al compilar};
\end{tikzpicture}%
}
\end{frame}

% =====================================================
\begin{frame}{Zephyr vs.\ un RTOS ``solo kernel''}
\small
\centering
\begin{tabular}{p{3.6cm}p{4.6cm}p{4.6cm}}
\toprule
& \textbf{Kernel solo} (p.\ ej.\ FreeRTOS) & \textbf{Zephyr} \\
\midrule
Scheduler, colas, semáforos & Sí & Sí \\
Drivers de periféricos & Los del fabricante (HAL propio) & Propios, con \alert{API común} \\
Descripción del hardware & En el código (\texttt{\#define}, init) & \alert{Devicetree} \\
Configuración & \texttt{FreeRTOSConfig.h} & Kconfig (\texttt{prj.conf}) \\
Build y flasheo & El del IDE del fabricante & \texttt{west} + CMake, igual para todos \\
Red, BLE, USB, FS & Bibliotecas externas & Integradas \\
\bottomrule
\end{tabular}

\vspace{0.8em}
\footnotesize No es mejor ni peor: \textbf{es otro alcance}. Zephyr ocupa más memoria y exige más formación.
\end{frame}

% =====================================================
\section{Devicetree}

% =====================================================
\begin{frame}{Devicetree: describir, no programar}
\begin{itemize}
    \item Estructura de datos \textbf{jerárquica} que describe el hardware: buses, periféricos, pines, direcciones
    \item Viene de Open Firmware y del kernel Linux (ARM, 2011)
    \item En Linux se interpreta \emph{al arrancar}; en Zephyr se resuelve \alert{al compilar}
    \begin{itemize}
        \item No hay parser en el microcontrolador
        \item Todo termina en \textbf{macros de C} $\Rightarrow$ costo cero en tiempo de ejecución
        \item Los errores (pin inexistente, propiedad faltante) aparecen \textbf{al compilar}
    \end{itemize}
\end{itemize}
\end{frame}

% =====================================================
\begin{frame}[fragile]{Anatomía de un nodo}
\begin{columns}[T]
\column{0.52\textwidth}
\begin{lstlisting}[language=dts]
/ {
    aliases {
        accel0 = &mpu6050;
    };
};

&i2c1 {
    status = "okay";
    clock-frequency = <400000>;

    mpu6050: mpu6050@68 {
        compatible = "invensense,mpu6050";
        reg = <0x68>;
    };
};
\end{lstlisting}
\column{0.45\textwidth}
\footnotesize
\begin{description}[\texttt{compatible}]
    \item[\texttt{/}] nodo raíz
    \item[\texttt{\&i2c1}] referencia a un nodo ya definido (por su \emph{label})
    \item[\texttt{mpu6050:}] \emph{label}: nombre para referenciarlo
    \item[\texttt{@68}] \emph{unit address}: dirección en el bus
    \item[\texttt{compatible}] \alert{qué es} el dispositivo $\Rightarrow$ qué driver lo maneja
    \item[\texttt{reg}] dirección I\textsuperscript{2}C
    \item[\texttt{status}] habilitado o no
    \item[\texttt{aliases}] nombres \alert{genéricos} que usa la aplicación
\end{description}
\end{columns}
\end{frame}

% =====================================================
\begin{frame}{Capas: del silicio a la aplicación}
\centering
\resizebox{0.95\textwidth}{!}{%
\begin{tikzpicture}[
    f/.style={draw, rounded corners, minimum width=3.3cm, minimum height=1.15cm, align=center, font=\small},
    arr/.style={-{Latex[length=2mm]}, thick, accentblue}
]
\node[f, fill=gray!20] (soc) {\texttt{rp2040.dtsi}\\\footnotesize SoC: periféricos,\\\footnotesize direcciones};
\node[f, fill=orange!20, right=0.8cm of soc] (brd) {\texttt{rpi\_pico.dts}\\\footnotesize placa: LED, cristal,\\\footnotesize pines usados};
\node[f, fill=green!20, right=0.8cm of brd] (ovl) {\texttt{rpi\_pico.overlay}\\\footnotesize aplicación: lo que\\\footnotesize yo conecté};
\node[f, fill=blue!15, right=0.8cm of ovl] (out) {\texttt{zephyr.dts}\\\footnotesize árbol final};
\draw[arr] (soc) -- (brd);
\draw[arr] (brd) -- (ovl);
\draw[arr] (ovl) -- (out);
\node[below=0.25cm of soc, font=\footnotesize\itshape] {lo escribe el fabricante};
\node[below=0.25cm of brd, font=\footnotesize\itshape] {lo escribe quien hace la placa};
\node[below=0.25cm of ovl, font=\footnotesize\itshape, text=accentblue] {\textbf{lo escribimos nosotros}};
\node[below=0.25cm of out, font=\footnotesize\itshape] {lo genera el build};
\end{tikzpicture}%
}

\vspace{1em}
\small Cada capa \textbf{agrega o modifica} nodos de la anterior. La aplicación nunca toca las dos primeras.
\end{frame}

% =====================================================
\begin{frame}{Bindings: el contrato entre el árbol y el driver}
\begin{columns}[c]
\column{0.5\textwidth}
\begin{itemize}
    \item Archivo YAML por cada \texttt{compatible}
    \item Define qué propiedades existen, su tipo y si son obligatorias
    \item El build \textbf{valida} el árbol contra el binding
    \item Zephyr v4.5.0 trae \alert{más de 4400} bindings
\end{itemize}
\column{0.46\textwidth}
\centering
\begin{tikzpicture}[b/.style={draw, rounded corners, minimum width=3.6cm, minimum height=0.8cm, align=center, font=\footnotesize}, arr/.style={-{Latex}, thick, accentblue}]
\node[b, fill=green!15] (n) {nodo\\\texttt{compatible = "invensense,mpu6050"}};
\node[b, fill=yellow!20, below=0.5cm of n] (y) {\texttt{invensense,mpu6050.yaml}};
\node[b, fill=blue!15, below=0.5cm of y] (d) {\texttt{drivers/sensor/tdk/mpu6050}};
\draw[arr] (n) -- node[right, font=\tiny] {se valida con} (y);
\draw[arr] (y) -- node[right, font=\tiny] {lo implementa} (d);
\end{tikzpicture}
\end{columns}
\end{frame}

% =====================================================
\begin{frame}[fragile]{Del árbol al código C}
\centering
\begin{tikzpicture}[
    f/.style={draw, rounded corners, minimum height=0.9cm, align=center, font=\footnotesize},
    arr/.style={-{Latex[length=2mm]}, thick, accentblue}
]
\node[f, fill=gray!15, minimum width=2.6cm] (in) {\texttt{.dtsi} + \texttt{.dts}\\+ overlays};
\node[f, fill=yellow!20, right=0.6cm of in, minimum width=2.4cm] (py) {scripts Python\\\footnotesize(edtlib + bindings)};
\node[f, fill=blue!15, right=0.6cm of py, minimum width=3.6cm] (h) {\texttt{devicetree\_generated.h}\\\footnotesize miles de \texttt{\#define}};
\node[f, fill=green!15, right=0.6cm of h, minimum width=2.4cm] (cc) {compilador\\de C};
\draw[arr] (in) -- (py); \draw[arr] (py) -- (h); \draw[arr] (h) -- (cc);
\end{tikzpicture}

\vspace{0.8em}
\begin{lstlisting}[language=C]
/* La aplicación pide "lo que sea que esté en led0"... */
static const struct gpio_dt_spec led = GPIO_DT_SPEC_GET(DT_ALIAS(led0), gpios);

/* ...y el preprocesador lo resuelve, por ejemplo en la Pico, a:
 *   { .port = &gpio0_device, .pin = 25, .dt_flags = GPIO_ACTIVE_HIGH }
 * Sin tablas en RAM, sin búsquedas en tiempo de ejecución. */
\end{lstlisting}
\end{frame}

% =====================================================
\begin{frame}[fragile]{Kconfig: la otra mitad de la configuración}
\begin{columns}[T]
\column{0.5\textwidth}
\textbf{\footnotesize prj.conf de la app del acelerómetro}
\begin{lstlisting}[language=kconf]
# APIs que usa la aplicación
CONFIG_GPIO=y
CONFIG_I2C=y
CONFIG_SENSOR=y

# Logging (con punto flotante)
CONFIG_LOG=y
CONFIG_CBPRINTF_FP_SUPPORT=y

# Shell para explorar el hardware
CONFIG_SHELL=y
CONFIG_I2C_SHELL=y
CONFIG_SENSOR_SHELL=y
\end{lstlisting}
\column{0.46\textwidth}
\begin{itemize}
    \item \textbf{Devicetree}: qué hardware hay
    \item \textbf{Kconfig}: qué software se compila
    \item Los drivers \alert{se activan solos}: si hay un nodo \texttt{"invensense,mpu6050"} con \texttt{status = "okay"}, se compila su driver
    \item Un mismo \texttt{prj.conf} para todas las placas
    \item \texttt{west build -t menuconfig} para explorar
\end{itemize}
\end{columns}
\end{frame}

% =====================================================
\section{Modelo de drivers}

% =====================================================
\begin{frame}[fragile]{\texttt{struct device}: todo dispositivo es igual}
\begin{columns}[T]
\column{0.5\textwidth}
\begin{lstlisting}[language=C]
struct device {
    const char *name;
    const void *config; /* flash */
    void *data;         /* RAM   */
    const void *api;    /* vtable */
    /* ... */
};
\end{lstlisting}
\footnotesize
\textbf{config}: datos del devicetree (constantes)\\
\textbf{data}: estado del driver\\
\textbf{api}: tabla de funciones del driver
\column{0.46\textwidth}
\begin{itemize}
    \item \textbf{Una instancia por nodo} del devicetree con \texttt{status = "okay"}
    \item Se crean \textbf{al compilar} (\texttt{DEVICE\_DT\_INST\_DEFINE})
    \item Se inicializan solas antes de \texttt{main()}, por niveles: \texttt{PRE\_KERNEL}, \texttt{POST\_KERNEL}\ldots
    \item La app obtiene el puntero con \texttt{DEVICE\_DT\_GET(nodo)}
\end{itemize}
\end{columns}
\end{frame}

% =====================================================
\begin{frame}{Una API, muchas implementaciones}
\centering
\resizebox{0.95\textwidth}{!}{%
\begin{tikzpicture}[
    b/.style={draw, rounded corners, minimum height=0.85cm, align=center, font=\small},
    arr/.style={-{Latex[length=2mm]}, thick, accentblue}
]
\node[b, fill=green!15, minimum width=6cm] (app) {\texttt{gpio\_pin\_toggle\_dt(\&led);}};
\node[b, fill=blue!10, below=0.6cm of app, minimum width=6cm] (api) {\texttt{gpio.h}: \texttt{api->port\_toggle\_bits(dev, BIT(pin))}};
\foreach \i/\txt/\c in {0/gpio\_rpi\_pico.c/orange, 1/gpio\_stm32.c/orange, 2/gpio\_esp32.c/orange, 3/gpio\_nrfx.c/orange, 4/gpio\_mcux\_lpc.c/orange}{
    \node[b, fill=\c!20, minimum width=2.6cm, font=\footnotesize\ttfamily] (d\i) at ($(api.south)+(-5.6+\i*2.8,-1.5)$) {\txt};
    \draw[arr] (api.south) -- (d\i.north);
}
\node[b, fill=gray!15, minimum width=14cm, below=0.6cm of d2] (hal) {Registros y HAL de cada fabricante};
\end{tikzpicture}%
}

\vspace{0.6em}
\small La función que se ejecuta la decide \textbf{el devicetree}, a través del puntero \texttt{api} del dispositivo.
\end{frame}

% =====================================================
\begin{frame}[fragile]{La API de sensores}
\begin{columns}[T]
\column{0.56\textwidth}
\begin{lstlisting}[language=C]
const struct device *s =
        DEVICE_DT_GET(DT_ALIAS(accel0));
struct sensor_value v[3];

sensor_sample_fetch(s);  /* medir */

sensor_channel_get(s,
        SENSOR_CHAN_ACCEL_XYZ, v);  /* m/s2 */
sensor_channel_get(s,
        SENSOR_CHAN_GYRO_XYZ, v);   /* rad/s */
/* canal que no mide -> -ENOTSUP */
\end{lstlisting}
\column{0.40\textwidth}
\begin{itemize}
    \item Canales con \textbf{unidades del SI}: m/s\textsuperscript{2}, rad/s, gauss, °C
    \item \texttt{sensor\_value}: entero + millonésimas, sin \texttt{float} obligatorio
    \item Igual para MPU6050, FXOS8700, FXLS8974, MMA8652FC y otros +240
    \item API asíncrona \texttt{sensor\_read()} + decodificadores para alto rendimiento
\end{itemize}
\end{columns}
\end{frame}

% =====================================================
\begin{frame}{APIs estandarizadas: no solo GPIO}
\centering
\small
\begin{tabular}{lll}
\toprule
\textbf{Periféricos} & \textbf{Comunicación} & \textbf{Servicios} \\
\midrule
GPIO, PWM, ADC, DAC & UART, I\textsuperscript{2}C, SPI, CAN & Logging, shell \\
Contadores, watchdog, RTC & USB device/host & Sistema de archivos \\
Sensores, LED, displays & Bluetooth LE & Settings (NVS) \\
Entrada (teclas, touch) & Ethernet, Wi-Fi, TCP/IP & Gestión de energía \\
Flash, EEPROM & Thread, Matter, LoRa & Actualización (MCUboot) \\
\bottomrule
\end{tabular}

\vspace{0.8em}
Mismo patrón en todas: \texttt{DEVICE\_DT\_GET} + \texttt{device\_is\_ready} + funciones de la API.
\end{frame}

% =====================================================
\section{Demo}

% =====================================================
\begin{frame}{El banco de pruebas: solo hardware integrado}
\centering
\footnotesize
\setlength{\tabcolsep}{4pt}\begin{tabular}{lllcclc}
\toprule
\textbf{Placa} & \textbf{Fabricante} & \textbf{Núcleo} & \textbf{LED} & \textbf{Botón} & \textbf{Acelerómetro} & \textbf{Temp.} \\
\midrule
nRF5340 DK & Nordic & Cortex-M33 & \si & \si & MPU6050 (ext.) & \no \\
FRDM-K64F & NXP & Cortex-M4F & \si & \si & FXOS8700 & \si \\
LPCXpresso55S69 & NXP & Cortex-M33 & \si & \si & MMA8652FC & \no \\
FRDM-MCXN236 & NXP & Cortex-M33 & \si & \si & FXLS8974 & \si \\
ESP32-C3 SuperMini & Espressif & \alert{RISC-V} & \si & \si & \no & \si \\
DOIT ESP32 & Espressif & \alert{Xtensa LX6} & \si & \si & \no & \no \\
Teensy 4.1 & PJRC (NXP) & Cortex-M7 & \si & \no & \no & \si \\
\bottomrule
\end{tabular}

\vspace{0.6em}
\footnotesize Cada demo usa lo que la placa ya trae. Única excepción: el MPU6050 del nRF5340 DK.\\
También compilan, como referencia de memoria: Pico, Pico 2, Black Pill F401 y ESP32-S3.
\end{frame}

% =====================================================
\begin{frame}[fragile]{App 1: \texttt{blinky\_boton} --- el \texttt{main.c} (lo esencial)}
\begin{lstlisting}[language=C]
static const struct gpio_dt_spec led   = GPIO_DT_SPEC_GET(DT_ALIAS(led0), gpios);
static const struct gpio_dt_spec boton = GPIO_DT_SPEC_GET(DT_ALIAS(sw0),  gpios);

int main(void)
{
    gpio_pin_configure_dt(&led, GPIO_OUTPUT_INACTIVE);
    gpio_pin_configure_dt(&boton, GPIO_INPUT);
    gpio_pin_interrupt_configure_dt(&boton, GPIO_INT_EDGE_TO_ACTIVE);
    gpio_init_callback(&boton_cb, boton_presionado, BIT(boton.pin));
    gpio_add_callback_dt(&boton, &boton_cb);

    while (1) {
        gpio_pin_toggle_dt(&led);
        k_msleep(periodos_ms[atomic_get(&indice)]);
    }
}
\end{lstlisting}
\centering\footnotesize Ni un número de pin. Ni un registro. Ni un \texttt{\#ifdef} de placa (solo ``¿hay botón?'').
\end{frame}

% =====================================================
\begin{frame}[fragile]{Lo que cambia por placa: nada}
\begin{columns}[T]
\column{0.5\textwidth}
\textbf{\footnotesize El botón es opcional}
\begin{lstlisting}[language=C]
#if DT_NODE_EXISTS(DT_ALIAS(sw0))
/* la placa tiene botón: se usa */
static const struct gpio_dt_spec boton =
    GPIO_DT_SPEC_GET(DT_ALIAS(sw0), gpios);
#else
/* sin botón: periodo fijo */
#endif
\end{lstlisting}
\footnotesize La Teensy 4.1 no tiene botón: la app lo sabe \textbf{al compilar}.
\column{0.46\textwidth}
\textbf{\footnotesize Ningún overlay para blinky}
\begin{itemize}
    \item Todas las placas de la demo ya definen \texttt{led0} (y casi todas \texttt{sw0}) en su \texttt{.dts}
\end{itemize}
\vspace{0.3em}
\begin{lstlisting}[language=sh]
west build -p -b nrf5340dk/nrf5340/cpuapp \
    apps/blinky_boton
west build -p -b esp32c3_supermini \
    apps/blinky_boton
west build -p -b teensy41 apps/blinky_boton
\end{lstlisting}
\end{columns}
\end{frame}

% =====================================================
\begin{frame}[fragile]{App 2: \texttt{acelerometro} --- no sabe qué sensor tiene}
\begin{lstlisting}[language=C]
/* "accel0": el nombre estándar de Zephyr para el acelerómetro de la placa */
static const struct device *const sensor = DEVICE_DT_GET(DT_ALIAS(accel0));

static const struct magnitud magnitudes[] = {
    {SENSOR_CHAN_ACCEL_XYZ,    3, "Aceleracion", "m/s2"},
    {SENSOR_CHAN_GYRO_XYZ,     3, "Giro",        "rad/s"},
    {SENSOR_CHAN_MAGN_XYZ,     3, "Campo magn.", "G"},
    {SENSOR_CHAN_DIE_TEMP,     1, "Temp. chip",  "C"},
    {SENSOR_CHAN_AMBIENT_TEMP, 1, "Temperatura", "C"},
};

while (1) {
    sensor_sample_fetch(sensor);
    for (size_t i = 0; i < ARRAY_SIZE(magnitudes); i++) {
        mostrar(&magnitudes[i]);  /* si el driver responde -ENOTSUP, se omite */
    }
    k_msleep(2000);
}
\end{lstlisting}
\end{frame}

% =====================================================
\begin{frame}[fragile]{Placa + sensor: overlays que se combinan}
\begin{columns}[T]
\column{0.47\textwidth}
\textbf{\footnotesize frdm\_k64f.dts} \footnotesize(de Zephyr, no lo tocamos)
\begin{lstlisting}[language=dts]
aliases { accel0 = &fxos8700; };
\end{lstlisting}
\textbf{\footnotesize boards/frdm\_k64f.overlay} \footnotesize(se aplica solo)
\begin{lstlisting}[language=dts]
sensor_i2c: &i2c0 {   /* mi bus I2C */
    status = "okay";
};
\end{lstlisting}
\textbf{\footnotesize mpu6050.overlay} \footnotesize(se pide al compilar)
\begin{lstlisting}[language=dts]
&sensor_i2c {
    mpu6050: mpu6050@68 {
        compatible = "invensense,mpu6050";
        reg = <0x68>;
    };
};
/ { aliases { accel0 = &mpu6050; }; };
\end{lstlisting}
\column{0.5\textwidth}
\begin{lstlisting}[language=sh]
# acelerómetro integrado: sin overlay extra
west build -p -b frdm_k64f apps/acelerometro
west build -p -b frdm_mcxn236 apps/acelerometro

# MPU6050 externo (el nRF5340 DK no trae)
west build -p -b nrf5340dk/nrf5340/cpuapp \
  apps/acelerometro \
  -- -DEXTRA_DTC_OVERLAY_FILE=mpu6050.overlay
\end{lstlisting}
\footnotesize
El overlay extra se aplica \textbf{después} del de la placa: si la placa tenía sensor, el MPU6050 lo reemplaza.\\[0.3em]
$N$ placas $+$ $M$ sensores = $N+M$ archivos, no $N\times M$.
\end{columns}
\end{frame}

% =====================================================
\begin{frame}{Mismo \texttt{main.c}, distinto sensor}
\centering
\small
\begin{tabular}{llcccc}
\toprule
\textbf{Sensor} & \textbf{Dónde} & \textbf{Acel.} & \textbf{Giro} & \textbf{Magn.} & \textbf{Temp.} \\
\midrule
MPU6050 & externo (nRF5340 DK) & \si & \si & \no & \si \\
FXOS8700 & integrado en FRDM-K64F & \si & \no & \si & \no \\
FXLS8974 & integrado en FRDM-MCXN236 & \si & \no & \no & \si \\
MMA8652FC & integrado en LPC55S69 & \si & \no & \no & \no \\
\bottomrule
\end{tabular}

\vspace{1em}
\begin{itemize}
    \item La app pide todas las magnitudes; cada driver \textbf{responde lo que mide}
    \item \texttt{accel0} ya lo definen las placas: la app no necesita overlay
    \item FXOS8700 y MMA8652FC: dos chips, \alert{un solo driver} (\texttt{nxp,fxos8700})
    \item El modo (solo acelerómetro o acelerómetro + magnetómetro) lo fija el \textbf{Kconfig de cada placa}
\end{itemize}
\end{frame}

% =====================================================
\begin{frame}[fragile]{App 3: \texttt{temperatura} --- el sensor interno del chip}
\begin{lstlisting}[language=C]
/* "die-temp0": el nombre estándar de Zephyr para la temperatura del chip */
static const struct device *const sensor = DEVICE_DT_GET(DT_ALIAS(die_temp0));

while (1) {
    sensor_sample_fetch_chan(sensor, SENSOR_CHAN_DIE_TEMP);
    sensor_channel_get(sensor, SENSOR_CHAN_DIE_TEMP, &temp);
    LOG_INF("Temperatura del chip: %.1f C", sensor_value_to_double(&temp));
    k_msleep(2000);
}
\end{lstlisting}
\vspace{0.3em}
\centering\small Casi todos los microcontroladores tienen uno. \alert{Cada fabricante lo mide distinto.}
\end{frame}

% =====================================================
\begin{frame}{Mismo alias, distinto hardware}
\centering
\footnotesize
\begin{tabular}{lll}
\toprule
\textbf{SoC} & \textbf{Cómo mide} & \textbf{Qué hace el overlay} \\
\midrule
ESP32-C3, ESP32-S3 & periférico dedicado & habilitar el nodo \\
i.MX RT1062 (Teensy 4.1) & monitor dedicado (TEMPMON) & habilitar el nodo \\
RP2040 (Pico) & canal 4 del ADC & habilitar el nodo \\
RP2350 (Pico 2) & canal 4 del ADC & nada: ya viene habilitado \\
STM32F401 (Black Pill) & ADC1 canal 16 + calibración & habilitar + alias \\
Kinetis K64F & ADC1 canales 26 y 27 & solo el alias \\
MCXN236 & LPADC0 canales 26A/26B & \alert{configurar el canal} \\
\bottomrule
\end{tabular}

\vspace{0.8em}
\footnotesize La aplicación siempre pide lo mismo: \texttt{die-temp0} y \texttt{SENSOR\_CHAN\_DIE\_TEMP}.\\
El nRF5340 tiene sensor, pero en el núcleo de red: desde el de aplicación no está disponible.
\end{frame}

% =====================================================
\begin{frame}[standout]
Demo en vivo

\Large Tres apps, tres arquitecturas,\\ni un \texttt{main.c} modificado
\end{frame}

% =====================================================
\begin{frame}{Resultados: lo que cuesta cada app}
\centering
\scriptsize
\begin{tabular}{llccc}
\toprule
\textbf{Placa} & \textbf{Núcleo} & \textbf{blinky\_boton} & \textbf{acelerometro} & \textbf{temperatura} \\
\midrule
Raspberry Pi Pico & Cortex-M0+ & 25,8 / 6,0 & 92,5 / 17,2 & 36,0 / 6,2 \\
Raspberry Pi Pico 2 & Cortex-M33 & 26,6 / 6,1 & 90,0 / 17,6 & 35,3 / 6,2 \\
Black Pill F401 & Cortex-M4F & 27,6 / 6,6 & 94,3 / 18,3 & 41,9 / 7,0 \\
DOIT ESP32 & Xtensa LX6 & 144,9 / 52,2 & 188,9 / 66,1 & -- \\
ESP32-C3 SuperMini & RISC-V & 131,8 / 91,1 & 211,9 / 116,3 & 134,1 / 93,6 \\
ESP32-S3 DevKitC & Xtensa LX7 & -- & 148,3 / 121,9 & 135,3 / 98,4 \\
nRF5340 DK & Cortex-M33 & 31,9 / 6,6 & 94,2 / 17,8 & -- \\
LPCXpresso55S69 & Cortex-M33 & 30,1 / 6,4 & 94,5 / 17,6 & -- \\
FRDM-MCXN236 & Cortex-M33 & 38,1 / 6,5 & 100,8 / 18,0 & 54,4 / 7,2 \\
FRDM-K64F & Cortex-M4F & 27,5 / 6,4 & 94,6 / 18,2 & 41,4 / 6,9 \\
Teensy 4.1 & Cortex-M7 & 42,1 / 6,5 & 104,8 / 15,3 & 51,0 / 6,5 \\
\bottomrule
\end{tabular}

\vspace{0.4em}
\footnotesize KiB de flash / RAM según el enlazador (Zephyr 4.5.0). 32 compilaciones, \alert{ningún \texttt{main.c} modificado}.\\
\texttt{acelerometro} incluye la shell; ``--'': la placa no tiene ese hardware.
\end{frame}

% =====================================================
\begin{frame}[fragile]{Explorar el hardware desde la shell}
\begin{columns}[T]
\column{0.52\textwidth}
\begin{lstlisting}
uart:~$ device list
uart:~$ i2c scan i2c@40005400
     0  1  2  3  4  5  6  7  8  9  a  b  ...
60: -- -- -- -- -- -- -- -- 68 -- -- --
uart:~$ sensor get mpu6050@68
uart:~$ kernel thread list
\end{lstlisting}
\column{0.44\textwidth}
\begin{itemize}
    \item La shell está en el \texttt{prj.conf}: viene en la app, sin recompilar nada
    \item Comandos para GPIO, I\textsuperscript{2}C, sensores, threads, memoria
    \item Ideal para validar el cableado antes de escribir código
\end{itemize}
\end{columns}
\end{frame}

% =====================================================
\section{Herramientas y límites}

% =====================================================
\begin{frame}[fragile]{\texttt{west}: una sola herramienta para todo}
\begin{columns}[T]
\column{0.5\textwidth}
\begin{lstlisting}[language=sh]
west init -l cie2026-zephyr   # workspace
west update                   # Zephyr + HALs
west build -b <placa> <app>
west flash                    # grabar
west debug                    # depurar (GDB)
west build -t menuconfig      # Kconfig
\end{lstlisting}
\column{0.46\textwidth}
\begin{itemize}
    \item \texttt{west flash} elige el programador según la placa: \texttt{pyocd}, \texttt{esptool}, \texttt{nrfutil}, \texttt{linkserver}, \texttt{teensy}, J-Link\ldots
    \item El \textbf{manifest} (\texttt{west.yml}) fija versiones $\Rightarrow$ builds reproducibles
    \item Funciona igual en Linux, macOS y Windows
\end{itemize}
\end{columns}
\end{frame}

% =====================================================
\begin{frame}{Lo que Zephyr \emph{no} resuelve}
\begin{itemize}
    \item \textbf{Periféricos exclusivos} de un fabricante (PIO del RP2040, radio de Nordic, ULP del ESP32): portables solo en parte
    \item \textbf{Soporte desparejo}: no todos los drivers implementan todo (DMA, bajo consumo, \emph{async})
    \item \textbf{Costo en memoria}: más que bare-metal o un kernel solo (blinky $\approx$30 KiB; con logging y shell 90--105 KiB en Cortex-M)
    \item \textbf{Cobertura desigual}: el sensor de temperatura no está disponible en el nRF5340 (núcleo de aplicación), el ESP32 ni el LPC55S69
    \item \textbf{Curva de aprendizaje}: devicetree, Kconfig, CMake y west de entrada
    \item \textbf{Temporización y consumo} siguen dependiendo del silicio: portar no es lo mismo que validar
\end{itemize}
\vspace{0.5em}
\centering\small La portabilidad no es gratis: \alert{se paga una vez}, en la arquitectura, y no en cada cambio de placa.
\end{frame}

% =====================================================
\begin{frame}{Para llevarse}
\begin{enumerate}
    \item El \textbf{devicetree} separa \emph{qué hardware hay} de \emph{qué hace la aplicación}
    \item Las \textbf{APIs estandarizadas} hacen que el driver sea un detalle de implementación
    \item Cambiar de placa o de sensor = \textbf{otro overlay}, no otro \texttt{main.c}
    \item El código de la aplicación se vuelve \textbf{reutilizable, testeable y enseñable}
\end{enumerate}
\end{frame}

% =====================================================
\begin{frame}{Probalo vos}
\begin{columns}[c]
\column{0.6\textwidth}
\begin{itemize}
    \item Workspace completo: manifest, apps y overlays
    \item Guía paso a paso para tu computadora (Linux, macOS o Windows)
    \item Nota técnica con los resultados
\end{itemize}
\vspace{0.5em}
\footnotesize\url{\repourl}

\vspace{0.8em}
\normalsize\textbf{Laboratorio de Sistemas Embebidos}\\
\footnotesize UTN -- Facultad Regional Avellaneda
\column{0.35\textwidth}
\centering
\qrcode[height=3.6cm]{\repourl}
\end{columns}
\end{frame}

% =====================================================
\begin{frame}[standout]
¿Preguntas?
\end{frame}

\end{document}
